\section{Background and construct definition}
Emotion regulation concerns attempts to influence emotional experience, expression or trajectory \parencite{gross2015}. Music-related regulation is heterogeneous: the Music in Mood Regulation scale and its brief version distinguish strategies \parencite{saarikallio2008,saarikallio2012}; the Healthy-Unhealthy Music Scale addresses potentially beneficial and problematic music engagement in a youth-oriented context \parencite{saarikallio2015}. These measures establish precedents, not a basis for copying their items or assuming validity for adult AI-assisted creators.

A general AI-assisted psychoemotional regulation measure, ERP-IA-6, was published in 2026 and provides initial evidence of a dominant dimension in Peruvian adults \parencite{briones2026}. Novelty cannot rest on being the first measure of AI-assisted emotion regulation. The narrower contribution is a profile of functions of active music creation, including unsuccessful sessions and potentially competing functions. Added value beyond existing measures remains an empirical question. Update the literature search before item freeze; an absence of identified comparators is not evidence of priority.

The AI-assisted Music for Emotion Regulation (AIM-ER) framework distinguishes motives, creative operations, mechanisms and outcomes. AI-MRFS targets perceived functions: what a particular activity served to do, as interpreted by the creator. It does not measure initial goals, acoustic quality, legal authorship, objective behaviour, affect change, clinical symptoms or regulatory success. Functions may co-occur. Temporary redirection of attention need not imply avoidance, and repeated attention to a theme need not establish rumination or harm.
\begin{table}[ht]
\centering\small
\caption{Provisional content domains, not an established factor structure.}
\begin{tabularx}{\linewidth}{lX}
\toprule
Domain & Intended scope \\
\midrule
EXP & Expressing and externalising personally relevant feelings.\\
REF & Exploring emotional meaning and alternative interpretations.\\
MOD & Attempting to influence felt emotional intensity or quality.\\
DIV & Temporarily redirecting attention away from emotional concerns.\\
SOC & Preparing or imagining communication of feelings to others.\\
PER & Repeatedly returning attention to the same emotional theme.\\
AVO & Using creation to postpone engagement with a concern.\\
\bottomrule
\end{tabularx}
\end{table}
\section{Research questions and provisional hypotheses}
RQ1 asks which domains are relevant, comprehensive and understandable. RQ2 asks which ordinal latent structure is supported in independent samples. RQ3 asks whether scores show adequate precision and construct-specific associations. RQ4 asks whether measurement is comparable across musical experience groups.
\begin{description}
\item[H1] A correlated multidimensional model frozen after exploratory analysis fits independent data better than a unidimensional model; no seven-factor model is assumed to survive refinement.
\item[H2] Each retained domain has reliability compatible with research use (target omega at least .70; report uncertainty).
\item[H3] EXP relates positively to independently assessed episode expression intention, REF to recalled post-episode clarity conditional on initial clarity, and MOD to initial mood-oriented motivation.
\item[H4] PER relates positively to a separately assessed tendency to dwell on negative experiences, and AVO to a separately assessed avoidance tendency.
\item[H5] EXP relates more strongly to expression intention than to frequency of generator use.
\end{description}
H3--H5 are provisional associations, not causal claims. If a domain is not retained or a comparator cannot be authorised before confirmatory freeze, the corresponding claim is not testable and cannot be substituted after inspection.
